\question[4]
Die Sicherheit des RSA-Verfahrens beruht darauf, dass das Produkt $m = p \cdot q$ zweier großer Primzahlen zwar leicht zu berechnen, aber praktisch nicht wieder in seine Primfaktoren zu zerlegen ist.
\begin{parts}
\part Erkläre, was man unter einer \textit{Einwegfunktion} versteht, und nenne die Einwegfunktion beim RSA-Verfahren.
\part Das Zerlegen einer Zahl mit 16 Stellen dauert auf einem Rechner etwa 1 Sekunde. Mit je zwei Stellen mehr wird die Rechenzeit zehnmal so lang. Wie lange dauert es ungefähr bei 26 Stellen?
\part Warum benutzt man in der Praxis Zahlen mit mehr als 600 Stellen?
\end{parts}
\begin{solutionbox}{4.5cm}
a) Eine Einwegfunktion lässt sich in einer Richtung leicht berechnen, die Umkehrung gelingt aber praktisch nicht in vertretbarer Zeit. Bei RSA ist das die Multiplikation zweier Primzahlen. \\
b) 10 Stellen mehr sind fünfmal der Faktor 10, also $10^5$ Sekunden, das ist mehr als ein Tag. \\
c) Für alle bekannten Angriffe wächst die Laufzeit so stark mit der Stellenzahl, dass ein Angriff bei so langen Zahlen praktisch nicht gelingt.
\end{solutionbox}
